\subsection{R 的紧子集}

定理 2（Borel-Lebesgue）. 实直线 R 的子集是紧的，其必要充分条件是它闭且有界。

\begin{enumerate}
\def\labelenumi{\arabic{enumi})}
\item
  条件是必要的。设 A 是 R 的紧子集，\(a\) 是一个 \(>0\) 的实数。集 A 是闭的（第一章 § 9 第 3 目，命题 4），且存在 R 的有限多个点 \(x_{i}\) （ \(1 \leqslant i \leqslant n\) ）使得 A 含于诸邻域 \([x_{i} - a, x_{i} + a]\) 的并（第一章 § 9 第 3 目）。令 \(b\) 为诸数 \(|x_{i}|\) 的最大者；则我们有 \(A \subset [-b - a, b + a]\) 。
\item
  条件是充分的。只需证明每个区间 \([-a, +a]\) （a \textgreater{} 0）是紧的；又由于这个区间是完备一致空间的一个闭子集，只需证明：对每个 b \textgreater{} 0 ，\([-a, +a]\) 都能用有限多个形如 \([x - b, x + b]\) 的区间覆盖（第二章 § 4 第 2 目，定理 3 的推论）。现在，设 n 是一个 \textgreater{} 0 的整数，使得 a \textless{} nb ；若 \(x \in [-a, +a]\) 且 m 是使 \(mb \leqslant x\) 的最大整数（正的或负的），则我们有 \(-n \leqslant m \leqslant n\) 且 \(mb \leqslant x \leqslant (m + 1)b\) 。于是 \(2n + 1\) 个区间 \([(k - 1)b, (k + 1)b]\) （ \(-n \leqslant k \leqslant n\) ）组成所要求类型的覆盖。
\end{enumerate}

推论 1. 实直线 R 的子集是相对紧的，当且仅当它有界。

推论 2. 实直线是局部紧空间且不是紧的。

注. 定理 2 常称为“Heine-Borel 定理”；参看第二章与第四章的历史注记。

